\section{Introduction}
\label{sec:intro}

\noindent The Great Depression led to a remarkable contraction in private investment that persisted throughout the first years of recovery: despite aggregate output growth of 11\% in 1933 and 9\% in 1934, net private investment remained negative in both years.\footnote{\cite{Cole2004} document that among common macroeconomic aggregates, investment was hit the hardest. Compared with gross national product (GNP), consumption, and hours worked, which were 36\%, 28\%, and 31\%, respectively, below trend in 1934, investment was 72\% below trend in the same year.} Given the lasting influence of the Great Depression on macroeconomics and finance, understanding the root causes behind the weak recovery of corporate investment is crucial, but the lack of disaggregated data over the period has limited researchers' ability to study this episode. Using newly hand-collected firm-level data, we show that U.S. public firms with greater exposure to the 1933 gold clause abrogation---which threatened widespread corporate bankruptcy---cut investment more severely during 1933 and 1934.

Prior to 1933, corporate bond issuers commonly indexed coupon and principal payments to the price of gold. In the era of the gold standard, these contractual agreements, known as the ``gold clauses'', aimed to protect bondholders against the devaluation of the dollar with respect to gold. When the United States effectively abandoned the gold standard in April 1933, the dollar quickly depreciated against gold, and the existence of these clauses posed a serious threat to most corporate bond issuers: around 97\% of outstanding corporate bonds were subject to a gold clause. Therefore, enforcing these clauses would have substantially increased the debt burden of nearly all corporate bond issuers during the depths of the worst economic collapse in recorded history.

To protect issuers against this threat, Congress passed in June 1933 a joint resolution that voided gold clauses in all existing and future contracts. Opposing the resolution, several bondholders subsequently filed lawsuits demanding payment in gold. Throughout the rest of 1933 and 1934, the validity of Congress's action faced litigation, with multiple lawsuits advancing through the judicial system. The litigation caused uncertainty for most corporate bond issuers about their future leverage. If the courts overturned the resolution, issuers would have to adjust bond payments according to the new gold price, increasing their leverage.

The Supreme Court delivered its ruling on four related court cases in February 1935. In a narrow 5-4 decision, the Court ruled that the abrogation of the gold clause in private contracts was constitutional; thus corporate bond payments were \textit{not} indexed to gold. Based on the price of gold at the time, had the Supreme Court instead overturned the abrogation, the coupon and principal payments on most corporate bonds would have increased by 69\%. Thus, while there were ultimately no changes in firm leverage as a result of the litigation, corporate bond issuers nonetheless experienced significant uncertainty over their future leverage during the episode.\footnote{In Section \ref{sec:historical_background}, we elaborate on the historical background relevant to our study.}
	
Guided by the dynamic debt overhang framework of \cite{Hennessy2004}, we test whether leverage uncertainty during 1933 and 1934 explains why firms with high exposure to gold clauses followed different investment paths than firms with low exposure. For this analysis, we hand-collect balance sheet and income statement data for all publicly traded U.S. industrial firms from 1925 to 1940. To measure each firm's exposure to the abrogation, we hand-collected data on the presence of gold clauses in individual securities---including corporate bonds and preferred equities---for all firms in our sample. Our hand-collected data reveal that although nearly all corporate bonds included such clauses, other types of long-term obligations, such as preferred shares, did not. Using this variation, we calculate the fraction of each firm's long-term liabilities that contained a gold clause.

Exploiting variation in gold clause exposure, we find that firms with fully exposed long-term liabilities had investment rates 6.0 and 4.1 percentage points lower in 1933 and 1934, respectively, than similar unexposed firms. Importantly, the negative effect of gold clause exposure is confined to the leverage uncertainty period of 1933--1934 and does not appear in any other year from 1926 to 1940. Moreover, examining equity payouts, we find that exposed firms increased their dividend payments during the leverage uncertainty period. 

These patterns are consistent with debt overhang as the mechanism driving the investment decline \citep{Myers1977, Hennessy2004}. The legal uncertainty over gold clauses raised the prospect of substantially higher liabilities for exposed firms, reducing equityholders' incentives to invest since bondholders would capture a larger share of returns from new capital. On the payout side, our results align with \cite{Hennessy2004}'s dynamic debt overhang framework, in which equityholders optimally extract large dividends from distressed firms.\footnote{See Section I.C in \cite{Hennessy2004}.} Consistent with debt overhang theory, we also find some evidence that both investment and payout effects are more pronounced among firms with lower credit ratings.

Our results complement the empirical literature on debt overhang that exploits plausibly exogenous variation in leverage to study firm decisions \citep[e.g.,][]{Giroud2012, Wittry2018, Bennett2025}. With respect to this literature, our empirical setting has a unique feature: the abrogation created uncertainty about future leverage increases that ultimately did not materialize. This characteristic connects our work with the literature on political and policy uncertainty \citep{Bloom2009, Baker2016, Hassan2019}. Our study is closest to the literature on firm-level political uncertainty pioneered by \cite{Hassan2019} as the abrogation of the gold clause represents political intervention in private debt contracts (\cite{Bolton2002}). We contribute to this literature by showing how an uncertainty shock to future liabilities delayed the recovery of investment during the Great Depression.

Evaluating the effects of gold clause uncertainty during the Great Depression presents significant empirical challenges given the numerous concurrent policy interventions and economic shocks of this period. A central component of our analysis therefore addresses potential confounding factors that might otherwise explain our findings.

A key concern is that firms' gold clause exposure was largely determined by their long-term liability composition (corporate bonds versus preferred equity), which is endogenously chosen rather than randomly assigned. This raises the possibility that firms strategically adjusted their gold clause exposure in anticipation of the abrogation. Such strategic adjustment could have been prompted by Britain's abandonment of the gold standard in September 1931, which created fears that the United States might follow suit and thus provided both warning and motivation for strategic adjustment.

To mitigate this concern, we provide empirical evidence that firms did not systematically adjust their gold clause exposure in the aftermath of Britain's departure from the gold standard. With the benefit of hindsight, this inertia can seem puzzling. However, foreign exchange data show that markets did not expect the United States to abandon the gold standard: even a few days before the suspension of convertibility, the dollar remained stable against currencies of countries that remained on the gold standard, suggesting widespread confidence that the United States would maintain its commitment to gold. This market expectation, combined with the substantial costs and constraints of restructuring long-term debt obligations, likely explains firms' apparent inaction. Nonetheless, to further mitigate endogeneity concerns, we measure firms' exposure to the gold clause in 1930, before Britain's departure from the gold standard.

Another challenge to our empirical approach is that 1933 and 1934 witnessed major policy shifts under the New Deal, including industrial reforms that could confound our results. For instance, the industrial codes implemented under the National Industrial Recovery Act of 1933 (NIRA) varied by industry and may have affected investment decisions differently across industries (\cite{Cole2004}). We show that the lower investment rates of high gold clause exposure firms are robust to industry-year fixed effects, which control for industry-specific policy effects. Beyond controlling for industry-year fixed effects, we also control for a battery of firm characteristics interacted with time. Overall, our results are robust across the various specifications, lending additional support to the conclusion that gold clause exposure, rather than other firm characteristics or New Deal policies, drove investment differences during this period.

While our investment and equity payout results are consistent with debt overhang predictions, firms likely faced other frictions during this period that could challenge our interpretation. For example, \cite{Benmelech2018} show that firms with maturing corporate bonds during the Great Depression experienced greater employment declines due to refinancing constraints. To ensure that refinancing constraints do not drive our investment results, we show that our findings hold using a subsample of firms without bonds maturing between 1931 and 1934. Our findings thus complement \cite{Benmelech2018} by showing that the gold clause embedded in bond contracts, in addition to rollover concerns, played a significant role in shaping firm outcomes during the Depression.

Beyond refinancing constraints on maturing bonds, other sources of financial friction could explain the investment gap we document. For example, the banking sector experienced severe disruptions during the Great Depression \citep{Bernanke1983}. However, in our sample of large industrial firms, bank debt accounts for less than 2\% of total liabilities, consistent with the view that larger firms are less affected by disruptions in the banking sector \citep{Bernanke2018, Siemer2014}. More broadly, we show that the investment gap in 1933 and 1934 is unique to corporate bond financing; exposure to other types of long-term liabilities that did not include the gold clause, such as preferred equity, show no similar effect on investment. These results further support the role of gold clause exposure in shaping firm outcomes during the Depression.

Finally, to quantify the importance of the gold saga at the aggregate level, we perform a simple partial equilibrium aggregation exercise based on our reduced-form estimates.\footnote{\cite{Chodorow2013, Hausman2017, Benmelech2018} present similar aggregation exercises.} Within our sample, we find that gold clause exposure accounts for one-third of public firms' divestment in 1933 and around one-half in 1934. This suggests that the abrogation of gold clauses played a substantial role in shaping corporate investment patterns of large industrial firms during a critical period of the Great Depression.

The outline of this article is as follows. Section \ref{sec:literature} summarizes the related literature. Section \ref{sec:historical_background} discusses historical events relevant to our study. Section \ref{sec:data} describes the firm-level data used in our analysis. Section \ref{sec:main_results} presents the empirical evidence on investment and the supporting evidence for the impact of gold clauses on firm outcomes. Section \ref{sec:conc} concludes.